\section{Background}

This section defines the problem setting (Sections~\ref{sec:efg}--\ref{sec:domains}), the
foundations our method builds on (Section~\ref{sec:cfr}), and the existing approaches to
speeding up CFR (Section~\ref{sec:variants}).

\subsection{Extensive-Form Games and Nash Equilibrium}
\label{sec:efg}

We model the problem as a finite \emph{extensive-form game} with imperfect
information~\cite{shoham2008multiagent}. Players in a set $N$ take turns acting, with a special
\emph{chance} player representing random events such as card deals. A \emph{history} $h$ is a
sequence of actions from the start of the game, $Z$ is the set of terminal histories, and
$u_i(z)$ is player $i$'s payoff at $z \in Z$. Histories that player $i$ cannot distinguish,
because they differ only in information hidden from $i$, are grouped into an \emph{information
set} $I \in \mathcal{I}_i$. The actions $A(I)$ available are the same at every history in $I$.

A \emph{strategy} $\sigma_i$ assigns a probability distribution $\sigma_i(I,\cdot)$ over $A(I)$
to every information set of player $i$, and a \emph{strategy profile} $\sigma$ contains one
strategy per player. We write $\pi^\sigma(h)$ for the probability of reaching $h$ under
$\sigma$, and $\pi^\sigma_{-i}(h)$ for the part of that probability contributed by chance and
all players other than $i$. The expected payoff is $u_i(\sigma) = \sum_{z\in Z}\pi^\sigma(z)\,
u_i(z)$.

A profile is a \emph{Nash equilibrium} if no player can improve its payoff by unilaterally
changing its strategy. In two-player zero-sum games ($u_1 = -u_2$), we measure distance from
equilibrium with \emph{exploitability},
\begin{equation}
  e(\sigma) = \frac{1}{2}\sum_{i\in\{1,2\}}\Big(\max_{\sigma_i'} u_i(\sigma_i',\sigma_{-i})
  - u_i(\sigma)\Big),
  \label{eq:expl}
\end{equation}
the average amount that a best-responding opponent could gain. Exploitability is zero exactly
at a Nash equilibrium. Exploitability is computed exactly by a best-response tree traversal, and is the
standard metric for comparing CFR variants~\cite{tammelin2015solving, brown2019solving}.

\subsection{Problem Domains}
\label{sec:domains}

We evaluate on two standard poker benchmarks small enough for exploitability to be computed
exactly.

\emph{Kuhn poker}~\cite{kuhn1950simplified} uses a three-card deck (Jack, Queen, King). Each of
two players antes one chip and receives one private card, followed by a single betting round
in which players may check, bet one chip, call or fold. If no one folds, the higher card wins. The game has 12 information sets, and its equilibria are known in
closed form, which makes it a useful correctness check.

\emph{Leduc poker}~\cite{southey2005bayes} uses a six-card deck containing two suits of Jack,
Queen and King. Each player antes one chip and receives one private card. After a first betting
round with bets of two chips, one public card is revealed, followed by a second round with bets
of four chips, with at most two raises per round. A player whose private card pairs the public
card wins; otherwise the higher card wins. Leduc poker adds two features that make full
Texas hold'em hard, multiple betting rounds and shared public information, while having only
936 information sets.

Both games are provided by OpenSpiel~\cite{lanctot2019openspiel}, an open-source framework that
implements the game rules, tabular policies, CFR and \cfrp{}, and exact exploitability.
OpenSpiel also provides three-player versions of both games.

\subsection{Regret Minimization and CFR}
\label{sec:cfr}

\emph{Regret} measures how much better an agent would have done had it always played a fixed
action instead of its actual sequence of strategies. An algorithm is \emph{no-regret} if its
average regret goes to zero. A central result connects learning to equilibrium: if both players
of a two-player zero-sum game have average regret at most $\epsilon$ after $T$ iterations of
self-play, their \emph{average} strategies form a $2\epsilon$-Nash
equilibrium~\cite{zinkevich2007regret}. \emph{Regret matching}~\cite{hart2000simple} is a
simple no-regret algorithm that plays each action in proportion to its positive cumulative
regret $R^T(a)$:
\begin{equation}
  \sigma^{T+1}(a) = \frac{\pos{R^T(a)}}{\sum_{b}\pos{R^T(b)}},
  \label{eq:rm}
\end{equation}
choosing uniformly at random if no regret is positive.

\emph{Counterfactual Regret Minimization}~\cite{zinkevich2007regret} applies regret matching
independently at every information set. The \emph{counterfactual value} of information set
$I \in \mathcal{I}_i$ is player $i$'s expected payoff from $I$, weighted by the probability
that chance and the opponents reach it:
\begin{equation}
  v_i^\sigma(I) = \sum_{h\in I}\sum_{z\in Z,\,h\sqsubseteq z}
  \pi^\sigma_{-i}(h)\,\pi^\sigma(h,z)\,u_i(z),
\end{equation}
where $\pi^\sigma(h,z)$ is the probability of reaching $z$ from $h$. Writing $v_i^\sigma(I,a)$
for the same value when $i$ plays $a$ at $I$, the \emph{instantaneous counterfactual regret}
at iteration $t$ is
\begin{equation}
  r^t_I(a) = v_i^{\sigma^t}(I,a) - v_i^{\sigma^t}(I).
  \label{eq:instregret}
\end{equation}
CFR accumulates $R^t_I(a) = R^{t-1}_I(a) + r^t_I(a)$ and picks the next strategy at each
information set with~\eqref{eq:rm}. Zinkevich et al.\ showed that the overall regret is bounded
by the sum of these local regrets, so the average strategy converges to a Nash equilibrium at a
rate of $O(1/\sqrt{T})$ in two-player zero-sum games.

In \emph{predictive} no-regret learning, the learner acts on a guess of its next loss, so its
regret depends on how wrong its guesses are rather than on the size of the losses, which can
speed up learning in games considerably~\cite{rakhlin2013optimization, syrgkanis2015fast}.

\subsection{Faster CFR Variants}
\label{sec:variants}

\textbf{\cfrp{}}~\cite{tammelin2015solving} makes three changes to CFR. First,
\emph{regret matching$^+$} keeps cumulative regrets non-negative,
\begin{equation}
  R^t_I(a) = \pos{R^{t-1}_I(a) + r^t_I(a)},
  \label{eq:rmplus}
\end{equation}
so an action that becomes good is played immediately. It also updates players in alternation
and weights iteration $t$ by $t$ in the average strategy, and was used to essentially solve
heads-up limit Texas hold'em~\cite{bowling2015heads}.

\textbf{Discounted CFR}~\cite{brown2019solving} multiplies accumulated positive regrets by
$t^\alpha/(t^\alpha+1)$, negative regrets by $t^\beta/(t^\beta+1)$, and average-strategy
contributions by $(t/(t+1))^\gamma$, with $\alpha=3/2$, $\beta=0$, $\gamma=2$ consistently
outperforming \cfrp{} in their experiments. These factors depend only on $t$ and are identical at every information
set.

\textbf{Predictive \cfrp{}}~\cite{farina2021faster} keeps update~\eqref{eq:rmplus} but plays
from the cumulative regret plus a prediction of the next instantaneous regret,
\begin{equation}
  \sigma^{t+1}_I(a) \propto \pos{R^t_I(a) + r^t_I(a)},
  \label{eq:pcfr}
\end{equation}
and \textbf{PDCFR$^+$}~\cite{xu2024minimizing} combines this prediction with DCFR-style
discounting.

Recent methods make these choices adaptive. \textbf{DDCFR}~\cite{xu2024dynamic} adjusts DCFR's
discount parameters during a run with a policy trained by reinforcement learning, while
\textbf{hyperparameter schedules}~\cite{zhang2026faster} change them as fixed functions of $t$.
\textbf{APCFR$^+$}~\cite{meng2025faster} scales the prediction in~\eqref{eq:pcfr} separately at
each information set, trusting it less where past predictions were inaccurate. \textbf{VAD-CFR},
found by an LLM-driven search over CFR variants~\cite{li2026discovering}, adjusts the discount
using a moving average of regret magnitude, although its authors found this component mattered
little on unseen games.

As Table~\ref{tab:prior} shows, the prediction has been adapted per information set, but the
discount on past regret only globally; \ours{} addresses this gap.

\begin{table}[t]
\centering
\caption{What existing CFR variants vary during a run, how, and at what scope.}
\label{tab:prior}
\setlength{\tabcolsep}{3pt}
\begin{tabular}{@{}llll@{}}
\toprule
\textbf{Method} & \textbf{Varies} & \textbf{How} & \textbf{Scope} \\
\midrule
\cfrp{}~\cite{tammelin2015solving}         & ---        & fixed         & --- \\
DCFR~\cite{brown2019solving}               & discount   & fixed in $t$  & global \\
\pcfrp{}~\cite{farina2021faster}           & ---        & fixed         & --- \\
PDCFR$^+$~\cite{xu2024minimizing}          & discount   & fixed in $t$  & global \\
DDCFR~\cite{xu2024dynamic}                 & discount   & learned (RL)  & global \\
HS~\cite{zhang2026faster}                  & discount   & schedule      & global \\
APCFR$^+$~\cite{meng2025faster}            & prediction & online        & per info.\ set \\
VAD-CFR~\cite{li2026discovering}           & discount   & online        & global \\
\midrule
\textbf{\ours{} (ours)}                    & \textbf{discount} & \textbf{online} & \textbf{per info.\ set} \\
\bottomrule
\end{tabular}
\end{table}
